\begin{textsample}{NEG Dim 1 – 2015\_c6 – Score -1.00 – 2015\_c6/enade\_2015\_ri\_4d06c54734.txt}  \label{ex:f1_neg_287}
CAPÍTULO 3 ANÁLISE TÉCNICA DA PROVA Este capítulo tem por objetivo apresentar o desempenho dos estudantes concluintes de Relações Internacionais no Enade/2015.
Para isso, foram calculadas as estatísticas básicas da prova em seu todo ( seção 3.1.1 ), bem como as estatísticas relacionadas aos Componentes de Formação Geral ( seção 3.1.2 ) e de Conhecimento Específico da Área ( seção 3.1.3 ).
Dadas as suas características, foram analisadas, em separado, as questões objetivas ( seção 3.2 ) e as questões discursivas ( seção 3.3 ).
Para as questões objetivas, foram disponibilizados os índices de facilidade e de discriminação ponto-bisserial, também em separado, para os Componentes de Formação Geral ( seção 3.2.1 ) e de Conhecimento Específico ( seção 3.2.2 ).
De \textbf{cada} componente, uma das questões foi escolhida para exemplificar a análise gráfica, relacionando as alternativas escolhidas pelos estudantes ( inclusive o gabarito ) com o número de acertos no componente.
O Anexo I apresenta a íntegra da análise gráfica para todas as questões objetivas.
Para \textbf{cada} uma das questões discursivas, os conteúdos dos tipos mais comuns de respostas dos estudantes são apresentados e comparados com o padrão de respostas esperado ( ver Anexo VIII com o padrão de respostas ).
Tomando como base as duas questões discursivas do Componente de Formação Geral, a seção 3.3.1.6 apresenta comentários sobre a correção das respostas com respeito à Língua Portuguesa.
Nas Tabelas 3.1 e 3.3 são apresentados o tamanho da população inscrita, ausente, presente e percentual de faltosos, respectivamente, por Grande Região e por Categoria Administrativa e Organização Acadêmica da IES.
As demais tabelas apresentam as seguintes estatísticas das notas15 : média do desempenho na prova, erro padrão da média, desvio padrão, nota mínima, mediana e nota máxima para \textbf{cada} um de seus componentes.
Tais estatísticas contemplam o total de estudantes concluintes da Área de Relações Internacionais inscritos e presentes à prova do Enade/2015, tendo em vista agregações ou por Grandes Regiões e o país como um todo, ou por Categoria Administrativa e Organização Acadêmica da IES.
Em relação aos gráficos de distribuição de notas, o intervalo considerado foi de 10 unidades, aberto à esquerda e fechado à direita, com exceção do primeiro intervalo, [ 0 ; 10 ], fechado em ambos os extremos.
Para os gráficos de distribuição das notas das 15 Uma definição dessas estatísticas pode ser encontrada no Glossário. questões discursivas, foram consideradas mais duas categorias : questão em branco16 e nota zero. 3.1 CONDIÇÃO DE PRESENÇA E ESTATÍSTICAS BÁSICAS DA PROVA Esta seção apresenta a condição de presença, além das estatísticas selecionadas e histogramas da nota geral ( 3.1.1 ) e de \textbf{cada} componente : Formação Geral ( 3.1.2 ) e Conhecimento Específico ( 3.1.3 ).
São também apresentadas estatísticas selecionadas de subpopulações caracterizadas por Grande Região, Categoria Administrativa e Organização Acadêmica. 3.1.1 Condição de Presença e Estatísticas Básicas Gerais A Tabela 3.1 apresenta a Condição de Presença na prova, por Grande Região, dos estudantes concluintes de Relações Internacionais.
A população total de Inscritos foi de 5.153.
Destes, 4.427 \textbf{estiveram} Presentes, sendo 14,1 \% o índice de não comparecimento.
A Região de maior abstenção foi a Centro-Oeste ( 21,3 \%, ) e a de menor abstenção foi a Sudeste ( 11,8 \% ).
Já a Tabela 3.2 apresenta as Estatísticas Básicas da prova, por Grande Região, dos estudantes concluintes de Relações Internacionais.
A Média das notas da prova como um todo ( nas seções seguintes serão analisados os Componentes de Formação Geral e de Conhecimento Específico ) foi 50,5, sendo que os alunos da região Norte obtiveram a Média mais baixa ( 45,6 ), e os da região Sul obtiveram a Média mais alta ( 50,9 ).
As demais médias foram : 50,8 na região Nordeste ; 50,7 na região Sudeste ; e 49,8 na região Centro-Oeste.
O Desvio padrão para o Brasil como um todo foi 14,7, sendo o maior Desvio padrão encontrado na região Sul ( 17,1 ) e o menor, na região Norte ( 11,6 ), indicando uma dispersão um pouco menor das notas desta última região.
A região que obteve a maior nota Máxima foi a Sul ( 92,8 ), ao passo que a região que atingiu a menor nota Máxima foi a Norte ( 78,1 ).
A Mediana do Brasil como um todo foi 50,8, sendo a maior Mediana obtida na região Sul ( 51,6 ), e a menor obtida no Norte ( 43,4 ).
A nota Mínima foi zero nas regiões Sudeste e Sul.
Considerando -se as notas segundo Grande Região, observa -se que existe diferença estatisticamente significativa ao nível de 95 \% entre a menor Média, obtida na região Norte ( 45,6 ), e as médias de todas as demais regiões.
O comportamento das notas dos estudantes de todo o Brasil pode ser observado no Gráfico 3.1 que apresenta um histograma com a distribuição das mesmas.
Essa é uma distribuição unimodal com a moda no intervalo ( 50 ; 60 ].
A Tabela 3.3 apresenta informações referentes à Condição de Presença desagregadas por Categoria Administrativa e por Organização Acadêmica.
Da população total de Inscritos, 3.684 são de IES Privada e 1.469 de IES Públicas.
Em relação à Organização Acadêmica, a maior participação foi obtida por estudantes de Universidades ( 3.204 ), seguido de Centros Universitários ( 1.064 ) e Faculdades ( 885 ).
A Categoria Administrativa de maior abstenção foi a Privada ( 15,4 \% ), e entre as Organizações Acadêmicas foi a das Universidades ( 14,2 \% ), que ficou acima da abstenção nacional de 14,1 \%.
Já a Tabela 3.4 apresenta informações referentes à Média da nota final desagregadas por Categoria Administrativa e por Organização Acadêmica.
A Média das notas da prova como um todo foi 50,5.
Em relação à Categoria Administrativa, os estudantes das IES Públicas obtiveram Média mais alta ( 55,6 ), e os das IES Privadas obtiveram Média mais baixa ( 48,3 ), que a Média nacional.
Observa -se que existe diferença estatisticamente significativa entre as médias das notas das IES Públicas e Privadas.
A diferença entre as médias das regiões Nordeste e Norte ( 4,8 ), a maior e a menor Média, é inferior à diferença entre IES Públicas e Privadas ( 7,3 ), caracterizando uma maior diversidade administrativa do que regional.
No tocante à Organização Acadêmica, apenas as Universidades obtiveram Média mais alta ( 51,6 ) que a nacional.
As demais médias foram : 47,9 nos Centros Universitários e 49,5 nas Faculdades.
O Desvio padrão para as IES Públicas ( 15,9 ) e para as Universidades ( 15,5 ) foi superior ao do Brasil como um todo ( 14,7 ), indicando uma dispersão um pouco maior das notas nesta Categoria Administrativa e nesta Organização Acadêmica.
Constata -se que existe diferença estatisticamente significativa ao nível de 95 \% nas médias das notas dos estudantes das Universidades em relação aos demais tipos de Organização Acadêmica. 3.1.2 Estatísticas Básicas no Componente de Formação Geral A Tabela 3.5 apresenta as Estatísticas Básicas em relação ao componente da prova que avalia a Formação Geral dos estudantes concluintes.
Os alunos de todo Brasil obtiveram desempenho médio de 65,5.
Quanto à variabilidade, o Desvio padrão das notas dos estudantes do Brasil, como um todo, foi 16,4.
A maior Média foi obtida na região Nordeste ( 66,7 ), e a menor, na região Norte ( 61,9 ).
As demais médias foram : 66,0 na região Sudeste ; 64,2 na região Sul ; e 65,2 na região Centro-Oeste.
Já o maior Desvio padrão foi obtido na região Sul ( 18,8 ), e o menor, na região Norte ( 14,8 ).
Os demais desvios padrões foram : 14,9 na região Nordeste, 15,7 da região Sudeste, e 16,3 na região Centro-Oeste.
A maior nota no Componente de Formação Geral da prova do Enade/2015 foi 99,0, obtida por, pelo menos, um aluno na região Centro-Oeste.
A menor nota Máxima foi obtida na região Sul ( 98,0 ).
A Mediana do Brasil, como um todo, foi 68,1, sendo a menor Mediana encontrada na região Norte ( 64,3 ) ; e a maior encontrada na região Nordeste ( 69,6 ).
A nota Mínima nesta parte foi zero em quase todas as regiões, exceto nas regiões Norte ( 29,5 ) e Nordeste ( 15,0 ).
Considerando -se as notas segundo Grande Região, observa -se que existe diferença estatisticamente significativa apenas entre as médias das notas do Componente de Formação Geral obtidas nas regiões Norte ( 61,9 ) e Sudeste ( 66,0 ).
O Gráfico 3.2 propicia a avaliação do desempenho dos estudantes do Componente de Formação Geral, a partir do histograma da distribuição das notas correspondentes.
A distribuição é unimodal, com moda em ( 70 ; 80 ], enquanto na prova, como um todo ( Gráfico 3.1 ), a moda foi alcançada dois intervalos abaixo ( 50 ; 60 ].
Nota -se, ainda que, no Gráfico 3.2, as notas apresentam maior dispersão do que no Gráfico 3.1 ( distribuição das notas da prova ), confirmado pela comparação dos desvios padrões : 14,7 para a nota da prova como um todo e 16,4 para o Componente de Formação Geral.
Na Tabela 3.6 são apresentadas as informações referentes ao desempenho dos concluintes do Componente de Formação Geral, em diferentes agregações : Categoria Administrativa e Organização Acadêmica.
Observa -se que existe diferença estatisticamente significativa entre as médias dos tipos de Categoria Administrativa.
A maior Média foi obtida por estudantes de IES Públicas ( 69,1 ), e a menor, obtida por estudantes de IES Privadas ( 64,0 ).
Considerando -se o tipo de Organização Acadêmica, nota -se que há diferença estatisticamente significativa entre as médias dos Centros Universitários em relação aos demais tipos de Organização Acadêmica. 3.1.3 Estatísticas Básicas do Componente de Conhecimento Específico A Tabela 3.7 apresenta as Estatísticas Básicas referentes ao Componente de Conhecimento Específico da área de Relações Internacionais.
A Média do desempenho dos alunos do Brasil, como um todo, foi 45,5.
A maior Média foi obtida na região Sul ( 46,5 ), e a menor, na região Norte ( 40,2 ).
As demais médias foram : 45,5 nas regiões Nordeste e Sudeste ; e 44,7 na região Centro-Oeste.
Quanto à variabilidade das notas, o Desvio padrão do Brasil, como um todo, foi 16,2, sendo o maior Desvio padrão observado na região Sul ( 18,2 ), e o menor, na região Norte ( 13,2 ).
Os demais desvios foram : 15,1 na região Nordeste, 15,5 na região Sudeste e 16,8 na região Centro-Oeste.
A Mediana das notas dos estudantes de todo o Brasil foi 45,2.
A maior Mediana ocorreu na região Sul ( 46,6 ), e a menor, na região Norte ( 39,5 ).
As demais medianas foram : 45,9 nas regiões Nordeste e Sudeste ; e 42,7 na região Centro-Oeste.
A nota Máxima do Brasil, como um todo, foi 95,5, sendo obtida por, pelo menos, um aluno na região Centro-Oeste.
As demais notas máximas foram : 76,1 na região Norte ; 78,6 na região Nordeste ; e 91,4 nas regiões Sudeste e Sul.
A nota Mínima foi zero nas regiões Sudeste e Sul.
Observa -se que existe diferença estatisticamente significativa da menor Média das notas do Componente de Conhecimento Específico da região Norte ( 40,2 ) em relação às demais regiões.
Assim como os Gráficos 3.1 e 3.2, o Gráfico 3.3, apresentado a seguir, permite uma avaliação do desempenho de concluintes de Relações Internacionais em relação ao Componente de Conhecimento Específico com um histograma da distribuição das notas correspondentes.
Dentre as três distribuições apresentadas, esta é a mais concentrada nas notas baixas.
Esta também é uma distribuição unimodal, e o grupo modal é o ( 40 ; 50 ], abaixo tanto do grupo modal da prova como um todo quanto do grupo modal para a Formação Geral.
A Tabela 3.8 apresenta uma comparação dos resultados em relação à Categoria Administrativa e à Organização Acadêmica, agora levando em conta o desempenho de alunos do Componente de Conhecimento Específico da prova.
No que se refere à Organização Acadêmica, a maior Média foi das Universidades ( 46,7 ), vindo a seguir a das Faculdades ( 44,1 ) e, depois, a dos Centros Universitários ( 43,0 ).
O maior Desvio padrão, e acima do valor para o Brasil, como um todo, foi o das Universidades ( 16,7 ).
As Universidades também obtiveram as maiores notas Máxima ( 95,5 ) e Mediana ( 47,4 ).
Por sua vez, os Centros Universitários obtiveram, respectivamente, 82,7 e 42,0 como nota Máxima e Mediana.
Já as Faculdades obtiveram, respectivamente, 89,9 e 44,1 como nota Máxima e Mediana.
A nota Mínima foi zero para Universidades, 3,7 para os Centros Universitários e 2,0 para as Faculdades.
Quanto à Categoria Administrativa, observa -se um comportamento semelhante àquele da parte de Formação Geral e à prova como um todo, ou seja, existe diferença 3.2 ANÁLISE DAS QUESTÕES OBJETIVAS Esta seção apresenta estatísticas selecionadas e histogramas das Questões Objetivas de Formação Geral ( 3.2.1 ).
São também apresentadas e comparadas as médias das subpopulações caracterizadas por Grande Região. 3.2.1 Componente de Formação Geral A Tabela 3.9 apresenta as Estatísticas Básicas relativas às oito questões objetivas do componente da prova que abrange a Formação Geral dos estudantes.
A Média do Brasil foi 66,1.
A menor Média foi encontrada na região Norte ( 58,9 ), e a maior, na região Nordeste ( 68,4 ).
As demais médias foram : 66,5 na região Sudeste ; 65,1 na região Sul ; e 66,2 na região Centro-Oeste.
O Desvio padrão do Brasil foi 20,2, sendo o maior Desvio padrão encontrado na região Sul ( 21,5 ), e o menor, na região Nordeste ( 18,6 ).
Os demais desvios foram : 20,5 na região Norte ; 19,9 na região Sudeste ; e 19,7 na região Centro-Oeste.
As medianas ( 62,5 ), as notas máximas ( 100,0 ) e as notas mínimas ( 0,0 ) foram iguais para quase todas as regiões, exceto a nota Mínima na região Norte ( 12,5 ) e a Mediana nas regiões Nordeste e Centro-Oeste ( 75,0 \textbf{cada} ). estatisticamente significativa entre as médias das IES Públicas ( 51,1 ) e IES Privadas ( 43,1 ).
Observa -se que existem diferenças estatisticamente significativas ao nível de 95 \% no Componente de Conhecimento Específico entre as notas das Universidades em relação aos Centros Universitários e Faculdades. 3.2 ANÁLISE DAS QUESTÕES OBJETIVAS Esta seção apresenta estatísticas selecionadas e histogramas das Questões Objetivas de Formação Geral ( 3.2.1 ).
São também apresentadas e comparadas as médias das subpopulações caracterizadas por Grande Região. 3.2.1 Componente de Formação Geral A Tabela 3.9 apresenta as Estatísticas Básicas relativas às oito questões objetivas do componente da prova que abrange a Formação Geral dos estudantes.
A Média do Brasil foi 66,1.
A menor Média foi encontrada na região Norte ( 58,9 ), e a maior, na região Nordeste ( 68,4 ).
As demais médias foram : 66,5 na região Sudeste ; 65,1 na região Sul ; e 66,2 na região Centro-Oeste.
O Desvio padrão do Brasil foi 20,2, sendo o maior Desvio padrão encontrado na região Sul ( 21,5 ), e o menor, na região Nordeste ( 18,6 ).
Os demais desvios foram : 20,5 na região Norte ; 19,9 na região Sudeste ; e 19,7 na região Centro-Oeste.
As medianas ( 62,5 ), as notas máximas ( 100,0 ) e as notas mínimas ( 0,0 ) foram iguais para quase todas as regiões, exceto a nota Mínima na região Norte ( 12,5 ) e a Mediana nas regiões Nordeste e Centro-Oeste ( 75,0 \textbf{cada} ).
A Tabela 3.10 apresenta o Índice de Facilidade e o Índice de Discriminação ( ponto-bisserial ) para \textbf{cada} uma das questões objetivas do Componente de Formação Geral.
Quanto ao Índice de Facilidade, foram usadas as seguintes cores para diferenciar o nível de dificuldade da questão :  Azul para as questões classificadas com índice Muito fácil ( \&gt ; =0,86 ), verde para as questões classificadas com índice Fácil ( 0,61 a 0,85 ), amarelo para as questões classificadas com Médio ( 0,41 a 0,60 ), vermelho para as questões classificadas com Difícil ( 0,16 a 0,40 ) e roxo para as questões classificadas com Muito difícil ( \&lt ; =0,15 ).
Já quanto ao Índice de Discriminação, foram usadas as seguintes cores para qualificar a questão :  As questões classificadas com índice Fraco receberam a cor vermelho ( \&lt ; =0,19 ), as classificadas com Médio receberam a cor amarelo ( 0,20 a 0,29 ), as classificadas com Bom receberam a cor verde ( 0,30 a 0,39 ) e as classificadas com Muito bom ( \&gt ; =0,40 ) receberam a cor azul.
As questões objetivas do Componente de Formação Geral, segundo o Índice de Facilidade, foram assim avaliadas : das oito questões, uma questão teve o Índice de Facilidade classificado como Muito fácil, situando -se na faixa de 0,86 ou mais ( 86,0 \% de acertos ou mais ).
Cinco questões foram tidas como Fácil, por terem índice de acertos situado na faixa entre 0,61 e 0,85 ( de 61,0 \% a 85,0 \% de acertos ).
Uma questão foi considerada de dificuldade Médio, situando -se no intervalo entre 0,41 e 0,60 do Índice de Facilidade, ou seja, houve entre 41,0 \% e 60,0 \% de acertos, enquanto uma questão foi classificada na categoria Difícil, situando -se no intervalo entre 0,16 e 0,40.
Por fim, nenhuma das questões apresentou menos de 15 \% de acertos, razão pela qual seria classificada como Muito difícil.
Como já comentado, para análise das questões objetivas, relativas à Formação Geral, segundo o poder de discriminação, utilizou -se o Índice de Discriminação ( ponto-bisserial ).
Nesta análise, as questões foram assim avaliadas : sete das oito questões apresentaram índice acima ou igual a 0,40 e, assim, foram classificadas com índice Muito bom para esse grupo de estudantes.
Uma questão teve Índice de Discriminação Bom, com valor entre 0,30 e 0,39, para esse grupo de estudantes.
Nenhuma questão teve nível Médio ou Fraco de discriminação para esse grupo de estudantes.
O Índice de Facilidade variou de 0,39 a 0,88, e o de Discriminação, de 0,38 a 0,55.
As sete questões com Índice de Discriminação Muito bom figuraram entre os diversos níveis de dificuldade desse conjunto : cinco classificadas na categoria Fácil ( questões 2, 3, 6, 7 e 8 ) do Índice de Facilidade, uma, na categoria Médio ( questão 4 ) e uma, na categoria Difícil ( questão 1 ).
Em particular, a questão 2 foi a que apresentou o maior poder discriminatório, com índice 0,55, porém foi considerada fácil em termos de facilidade, com uma proporção de 0,72 acertos.
O máximo de acertos foi alcançado pela questão 5 com um Índice de Facilidade de 0,88.
A questão de número 4 apresentou um Índice de Facilidade de 0,41, ou seja, um pouco menos da metade dos estudantes conseguiu resolvê -la, dentro do universo de participantes.
Seu Índice de Discriminação foi Muito bom ( 0,43 ).
Já a questão 1 obteve Índice de Discriminação Muito bom, 0,40, porém seu Índice de Facilidade foi Difícil ( 0,39 ).
O Gráfico 3.4, para exemplificar, analisa o comportamento da questão de número 2 de Formação Geral.
Trata -se de uma questão considerada mediana, com relação à facilidade e a que obteve o maior Índice de Discriminação dessa parte da prova.
Neste gráfico, \textbf{cada} uma das cinco curvas representa o percentual de respostas em determinada alternativa da questão, em função do número de acertos dos estudantes nessa parte da prova ( Formação Geral/Múltipla Escolha ), antes de possíveis eliminações pelo critério do ponto bisserial.
A curva em preto corresponde à alternativa C, a correta para esta questão.
Assim, observa -se que entre os estudantes com menor número de acertos, nessa parte do exame, a situação mais frequente foi a escolha da alternativa D ( em roxo ), incorreta.
Por exemplo, entre os estudantes que acertaram três questões, 35,2 \% escolheram a alternativa D, 30,8 \% a alternativa C ( correta ), 11,0 \% a alternativa A, 15,9 \% a alternativa B e 5,8 \% a alternativa E.
Entre os que acertaram três respostas entre as questões de múltipla escolha de formação geral, 0,6 \% deixou a questão em branco e 0,9 \% marcou mais de uma alternativa, invalidando a questão.
À medida em que o número de acertos aumenta, indicando desempenho melhor nessa parte da prova, aumenta concomitantemente a proporção de estudantes que selecionaram a alternativa correta C, atingindo 100 \% para os estudantes com 8 acertos.
Essa análise permite verificar como a questão discriminou os grupos de desempenho, justificando o alto índice obtido na questão.
Cumpre notar que não é possível inferir deste gráfico nem o índice de facilidade ( que seria uma média da proporção ponderada pela quantidade de alunos com \textbf{cada} uma das notas ), nem o índice de discriminação ponto bisserial, por razão equivalente.
No caso extremo no qual a grande concentração dos acertos dos alunos fosse abaixo de 4, o índice de facilidade seria obrigatoriamente abaixo de 40 \% ( neste exemplo ).
Caso a concentração fosse em 6 acertos ou mais, o índice seria obrigatoriamente acima de 60 \%.
Os gráficos relativos às demais questões de Formação Geral constam do Anexo I. 3.2.2 Componente de Conhecimento Específico A Tabela 3.11 apresenta as Estatísticas Básicas em relação às questões objetivas do Componente de Conhecimento Específico da prova de Relações Internacionais por Grande Região.
A Média do Brasil deste componente foi de 47,5.
A menor Média foi observada na região Norte ( 42,6 ), e a maior, na região Sul ( 48,7 ).
O Desvio padrão de todo o Brasil foi 16,8, sendo o menor Desvio padrão encontrado na região Norte ( 13,8 ), e o maior, na região Sul ( 18,9 ).
A Mediana de todo o Brasil foi 47,8, a mesma encontrada em todas as regiões do país, exceto na região Norte ( 43,5 ).
A nota Máxima da prova foi 100,0, obtida nas questões objetivas do Componente de Conhecimento Específico, por, pelo menos, um aluno na região Centro-Oeste, e nas demais regiões, a nota Máxima da prova foi : 78,3 nas regiões norte e Nordeste ; 95,7 na região Sudeste ; e 91,3 na região Sul.
Em quase todas as regiões a nota Mínima foi zero, exceto nas regiões Norte ( 17,4 ) e Centro-Oeste ( 8,7 ).
A Tabela 3.12 apresenta os Índices de Facilidade e Discriminação ( ponto-bisserial ) das questões objetivas do Componente de Conhecimento Específico para os estudantes de Relações Internacionais.
Para facilitar a diferenciação das questões, usaremos as mesmas cores da Tabela 3.10 para as diferentes classificações dos Índices de Facilidade e de Discriminação.
Dentre as questões objetivas da parte da prova relativa ao Componente de Conhecimento Específico, nenhuma questão foi anulada pela Comissão.
Desse modo, a classificação, quanto ao Índice de Facilidade, foi estabelecida com base em todas as 27 questões.
A partir dos índices obtidos, pode -se concluir que um pouco mais da metade das questões objetivas da prova foi considerado, pelo menos, Difícil : das 27 questões, 15 foram classificadas como Difícil ( a classificação modal ) ou como Muito difícil.
Cinco questões foram classificadas como Fácil ou como Muito fácil, e outras sete consideradas como Médio.
Já quanto aos índices de discriminação das questões objetivas do Componente de Conhecimento Específico da prova, tem -se como resultado a seguinte classificação : dez das 27 questões foram consideradas como boas, enquanto seis delas tiveram Índice de Discriminação Muito bom.
Assim, para a maioria das questões – 16 em 27 – os Índices de Discriminação foram Bom ou Muito bom.
Dentre as demais, sete delas foram classificadas como Médio e outras quatro como Fraco, sendo 11, por conseguinte, a quantidade de questões nos dois patamares mais baixos de discriminação.
Constata -se, assim, que a prova – no que se refere ao Componente de Conhecimento Específico – possuía boa capacidade de discriminar entre aqueles que dominam ou não o conteúdo.
Dentre as questões que alcançaram os maiores índices de discriminação, seis delas, as de números 10, 17, 18, 19, 25 e 32, foram classificadas com Índice de Discriminação Muito bom, situando -se no intervalo de 0,40 a 0,50 do índice, três delas ( questões 17, 19 e 25 ) foram classificadas na categoria Fácil, quanto ao Índice de Facilidade, e as outras três ( questões 10, 18 e 32 ), na categoria Médio.
A questão de número 34 foi a mais difícil dentre as 27 questões específicas, com baixo Índice de Facilidade, apenas 16,0 \% de acertos.
Essa questão apresentou poder discriminatório igualmente baixo, 0,03, o que comprova ter sido esta questão a mais difícil para os estudantes.
Destaca -se, também, a questão 9, com Índice de Facilidade 0,18, o que, em termos percentuais, corresponde a 18,0 \% de estudantes que responderam acertadamente.
Já 0,11 foi o de Índice de Discriminação.
Tais questões foram, portanto, pelo critério ponto-bisserial, consideradas inadequadas.
Por isso, as questões 9 e 34 foram eliminadas do cômputo da nota final.
Além destas duas, as demais questões com Índice de Discriminação Fraco, questões 16 e 31 também não foram utilizadas no cômputo final das notas, num total de quatro questões eliminadas.
A título de exemplo das análises do comportamento das questões objetivas, o Gráfico 3.5 analisa a questão 10 do Componente de Conhecimento Específico.
Apesar de esta questão ter sido considerada pelos estudantes avaliados como uma questão de classificação Médio como nível de facilidade da prova, apresentou Índice de Facilidade 0,53, ou seja, 53,0 \% dos estudantes assinalaram acertadamente a opção D, A título de exemplo das análises do comportamento das questões objetivas, o Gráfico 3.5 analisa a questão 10 do Componente de Conhecimento Específico.
Apesar de esta questão ter sido considerada pelos estudantes avaliados como uma questão de classificação Médio como nível de facilidade da prova, apresentou Índice de Facilidade 0,53, ou seja, 53,0 \% dos estudantes assinalaram acertadamente a opção D, correspondente ao gabarito.
Seu Índice de Discriminação foi igual a 0,50, classificado como Muito bom, também sendo esta questão a que apresentou o maior índice discriminatório.
Neste gráfico, \textbf{cada} uma das cinco curvas representa o percentual de respostas em determinada alternativa da questão 10, em função do número de acertos dos estudantes nessa parte da prova, antes de possíveis eliminações de questões pelo critério do ponto-bisserial.
A alternativa correta D, representada no gráfico pela curva em roxo, foi escolhida em maiores proporções pelos alunos com desempenho melhor nessa parte da prova.
Já as alternativas incorretas, também denominadas distratores, foram selecionadas, principalmente, por aqueles com notas mais baixas.
Neste caso também a soma não é sempre 100 \% por causa das questões não respondidas ou com mais de uma opção marcada.
Aqueles com nota zero, na sua quase totalidade deixaram esta questão em branco ou marcaram mais de uma alternativa, comportamento considerado inválido.
A proporção de alunos que selecionou a resposta correta D aumenta gradativamente, chegando a atingir 100 \% para 21 acertos ou mais, enquanto a proporção dos que escolheram alternativas incorretas decai, a partir de dois acertos, como função do número de acertos nessa parte da prova.
Os gráficos relativos às demais questões do Conhecimento Específico constam do Anexo I. 3.3 ANÁLISE DAS QUESTÕES DISCURSIVAS Esta seção apresenta estatísticas selecionadas e histogramas das Questões Discursivas de Formação Geral ( 3.3.1 ) e Conhecimento Específico ( 3.3.2 ).
São também apresentadas e comparadas as médias de subpopulações caracterizadas por Grande Região. 3.3.1 Componente de Formação Geral As análises dos resultados de desempenho dos estudantes de Relações Internacionais nas duas questões discursivas relativas à Formação Geral encontram -se na Tabela 3.13 e no Gráfico 3.6.
Na Tabela 3.13, observa -se que a nota Média nesse conjunto de questões foi próxima da obtida nas objetivas.
Os estudantes, de todo o Brasil, obtiveram, em Formação Geral, Média 66,1 nas questões objetivas e 64,6 nas questões discursivas.
No entanto, pode -se notar um aumento do Desvio padrão de 20,2 nas questões objetivas do Componente de Formação Geral dos alunos de todo o Brasil, para 22,1 nas questões discursivas do mesmo componente.
A maior Média foi obtida na região Norte ( 66,5 ), e a menor, na região Sul ( 62,7 ).
A Mediana de todo o Brasil, neste componente, foi 70,0, a mesma obtida nas regiões Norte e Nordeste.
Na região Sudeste ( 70,5 ) a Mediana foi maior, e nas regiões Sul e Centro-Oeste ( 69,5 \textbf{cada} ), foi menor.
A nota Máxima ( 99,0 ) foi obtida na região Sudeste ; sendo 96,5 nas regiões Norte e Nordeste, 98,0 na região Sul e 97,5 na região Centro-Oeste.
A nota Mínima ( 0,0 ) foram as mesmas em todas as regiões do Brasil, sem exceção.
O Gráfico 3.6 representa a distribuição das notas nas questões discursivas do Componente de Formação Geral.
A moda desta distribuição ocorre no intervalo ( 70 ; 80 ].
A Mediana de todo o Brasil, neste componente, foi 70,0, a mesma obtida nas regiões Norte e Nordeste.
Na região Sudeste ( 70,5 ) a Mediana foi maior, e nas regiões Sul e Centro-Oeste ( 69,5 \textbf{cada} ), foi menor.
A nota Máxima ( 99,0 ) foi obtida na região Sudeste ; sendo 96,5 nas regiões Norte e Nordeste, 98,0 na região Sul e 97,5 na região Centro-Oeste.
A nota Mínima ( 0,0 ) foram as mesmas em todas as regiões do Brasil, sem exceção.
O Gráfico 3.6 representa a distribuição das notas nas questões discursivas do Componente de Formação Geral.
A moda desta distribuição ocorre no intervalo ( 70 ; 80 ].
Na sequência, os resultados verificados para \textbf{cada} uma das questões discursivas de Formação Geral serão apresentados, estabelecendo -se relações com os conteúdos abordados em \textbf{cada} uma delas.
Os comentários da Banca de docentes corretores, a respeito do observado na correção das respostas dos estudantes, suas impressões e conclusões serão apresentados junto à análise de \textbf{cada} questão.
Cumpre esclarecer que, tendo em vista que as questões discursivas de Formação Geral são padronizadas, ou seja, constam de todas as provas, os comentários da Banca são os mesmos para todas as carreiras acadêmicas, sendo direcionados a todos os estudantes que participaram do Enade/2015.
A seguir, serão analisados os desempenhos da Área de Relações Internacionais nas duas questões discursivas de Formação Geral do Enade/2015, comparando -se os resultados obtidos com comentários para \textbf{cada} questão. 3.3.1.1 Análise de Conteúdo da Questão Discursiva 1 do Componente de Formação Geral Os dados de Relações Internacionais, obtidos a partir das respostas à questão 1, encontram -se na Tabela 3.14 e no Gráfico 3.7.
Nessa questão – de melhor desempenho dentre as duas de Formação Geral – os alunos, de todo o Brasil, tiveram Média 68,2.
A maior Média para a questão 1 foi obtida na região Norte ( 70,7 ), e a menor, na região Sul ( 65,5 ).
Quanto à variabilidade das notas, o Desvio padrão de todo o Brasil foi 23,0.
O menor Desvio padrão foi obtido na região Norte ( 16,9 ), e o maior Desvio padrão foi obtido na região Sul ( 26,1 ).
As medianas ( 75,0 ), as notas mínimas ( 0,0 ) e as notas máximas ( 100,0 ) foram as mesmas para todas as regiões do Brasil, sem exceção.
O Gráfico 3.7 mostra a distribuição das notas na questão discursiva 1 do Componente de Formação Geral.
Observa -se que a maior frequência corresponde aos alunos que obtiveram nota dentro do intervalo ( 70 ; 80 ], correspondendo à moda da distribuição. 3.3.1.2 Comentários sobre a correção de Conteúdo das respostas à Questão Discursiva 1 O enunciado da Discursiva 1 de Formação Geral \textbf{estava} claro e indicava um tema de extrema relevância, qual seja : a busca pela educação universal, especificamente, que apontava um caso onde ocorreu a discriminação de acesso ao ensino formal atingindo uma menina, representando a categoria das mulheres discriminadas em tantos campos da vida e, mais ainda, em sociedades que cultuam o poder soberano dos homens.
A linguagem utilizada foi clara, indicando a existência dos fatos vinculados à luta da menina Malala pelo acesso ao ensino formal das jovens do seu país, e o confronto ideológico que gerou o ataque à sua integridade física.
Diante da repercussão internacional, houve o reconhecimento pelo prêmio Nobel, e a ONU resolveu apoiá -la como símbolo da luta pela igualdade entre os sexos buscando a universalização da educação.
O conteúdo \textbf{está} integrado às Diretrizes, vez que o acesso ao ensino formal \textbf{está} vinculado ao direito constitucional à educação, que, também, é protegido pelo Estatuto da Criança e do Adolescente e nos pactos de Direitos Humanos aderidos pelo Brasil.
Por outro lado, o texto demandava uma reflexão sobre a realidade dos indivíduos que não têm acesso à educação e a situação das mulheres excluídas desse processo por segmentos da sociedade.
Como se tratava de texto contendo uma ampla gama de possibilidades, era natural a ocorrência de interpretações diversas das esperadas.
Releva notar que o enunciado pretendeu, primordialmente, identificar a questão de gênero, apontando para as dificuldades que as mulheres possuem, não somente no Brasil, mas em vários países do mundo.
Ao analisar as provas realizadas, constatamos que no item ‘ a ’ poucos foram os que identificaram os estatutos normativos indicados no padrão de resposta.
Essa ausência pode ser justificada tendo em vista que, no enunciado indicado, não existia elemento que pudesse ser utilizado pelo concluinte para embasar sua resposta nesse caminho.
As respostas foram mais gerais, identificando que a educação constitui um item fundamental para o desenvolvimento do indivíduo e que as mulheres sofrem mais com as limitações impostas por países com regime patriarcal, que discriminam o sexo feminino e, com frequência, proíbem as mulheres de frequentar a escola.
Deve -se ter em vista que a indicação de resposta que vinculava o texto aos estatutos normativos representou um raciocínio próprio dos estudantes de Direito que, naturalmente, encaminharam suas respostas para a fundamentação nessa linha identificando, não somente a Constituição, mas também leis, regulamentos, etc.
Tal linha de argumentação não foi objeto da formação dos estudantes das demais áreas.
A maior parte dos concluintes estabeleceu o tema educação formal associado à busca pelo conhecimento, como acesso ao ensino, indicando a independência intelectual e o desenvolvimento da consciência crítica, limitando as manipulações e permitindo amplo acesso ao mercado de trabalho, traduzindo uma libertação econômica e ascensão social.
No tocante ao item ‘ b ’, restou claro que a maioria esmagadora das respostas confluiu para as indicações do padrão apresentado, sendo relevante o tema de igualdade de gênero, além de candente para as concluintes do sexo feminino que desenvolveram os textos de forma concatenada, fundamentada e, muitas vezes, lamentando que, nos dias atuais, ainda se discuta a necessidade de igualar os gêneros na sociedade.
Nessa linha, foram identificadas respostas, consoante com o padrão apresentado, especialmente nos seguintes temas : • violência física e psicológica contra a mulher, incluindo a Lei Maria da Penha, no caso específico do Brasil ; • tolerância/intolerância a vestimentas, trajes, comportamentos socialmente estereotipados ; • aspectos socioculturais que impõem à mulher uma condição de submissão na sociedade, tais como a impossibilidade de a mulher manifestar seus desejos e posicionamentos em algumas culturas, entre outros ; • igualdade/desigualdade de gênero, por exemplo no mercado de trabalho, em relação à desigualdade salarial ; • ideais de Liberdade, Igualdade e Fraternidade : referência a esses ideais como possibilidade de equilibrar as relações de poder entre homens e mulheres.
Em muitos casos, também, os concluintes fizeram um liame entre o acesso à educação e a melhora das relações entre homens e mulheres ao permitir que a conclusão de maior grau de educação formal permita o acesso a cargos públicos relevantes, empregos bem remunerados e posições de poder, tanto na área pública quanto na área privada.
Como as opções do segundo item foram mais amplas, o padrão de resposta identificou as linhas de pensamento adotadas pelos concluintes.
A maior parte das respostas identificou, claramente, o problema da relação de gênero e suas consequências na sociedade.
Também foi frequente constatarem que existe efetiva melhora na situação das mulheres no Brasil, mas que tal movimento não ocorre em todo o mundo.
Aliás, o exemplo de Malala conduziu a uma reflexão em âmbito mundial, e os concluintes apresentaram suas análises críticas à situação da mulher em países que não admitem sequer a sua inserção no plano educacional e, muito menos, no mercado de trabalho.
A maior parte dos concluintes apresentou concatenação lógica nas suas respostas e uma minoria, pouco expressiva, não demonstrou esse domínio de 3.3.1.3 Análise de Conteúdo da Questão Discursiva 2 do Componente de Formação Geral A Tabela 3.15 mostra que o desempenho médio dos estudantes na questão discursiva 2 ( média 62,4 ) foi inferior ao obtido na questão discursiva 1 ( média 68,2 ).
A região Norte foi aquela cuja Média, nessa questão, foi maior ( 64,5 ), e a de menor Média foi a região Centro-Oeste ( 60,4 ).
Quanto à variabilidade das notas, o Desvio padrão de todo o Brasil foi 25,3, superior ao obtido na questão discursiva 1 ( 23,0 ).
O maior desvio nessa questão foi obtido na região Nordeste ( 27,8 ), enquanto o menor foi obtido na região Norte ( 21,4 ).
A Mediana de todo o Brasil foi 64,0, a mesma das regiões Nordeste, Sudeste e Centro-Oeste, enquanto a maior Mediana foi encontrada na região Norte ( 66,0 ) e a menor, na região Sul ( 63,0 ).
As notas máximas ( 100,0 ) foram obtidas nas regiões Nordeste e Sudeste.
As notas mínimas ( 0,0 ) foram as mesmas em todas as regiões do Brasil, sem exceção.
O Gráfico 3.8 mostra a distribuição das notas na questão discursiva 2 do Componente de Formação Geral.
Observa -se que a maior frequência corresponde aos alunos que obtiveram nota situada no intervalo ( 60 ; 70 ]. 3.3.1.4 Comentários sobre a correção de Conteúdo das respostas à Questão O enunciado da Discursiva 2 \textbf{estava}, igualmente, claro, indicando um tema de extrema relevância, qual seja : o reconhecimento de manifestações artísticas nascidas na periferia, como integrante da cultura reconhecida pela sociedade.
A linguagem utilizada foi clara, indicando a análise dos fatos e suas consequências no cotidiano dos partícipes e o confronto com outras visões culturais que são consideradas mais adequadas pela sociedade civil.
O conteúdo \textbf{está} integrado às Diretrizes, vez que o aspecto cultural e sua inter-relação com o conhecimento haurido no nível superior é essencial para a compreensão de mundo do universitário concluinte, com reflexo na sua atuação profissional.
O texto demanda uma reflexão sobre a realidade dos indivíduos que não têm acesso aos itens básicos da sociedade, como educação, saúde, cultura tradicional e segurança e que, apesar disso, criam alternativas para o seu lazer.
A incorporação do funk como manifestação cultural e artística da sociedade brasileira encontrou, basicamente, duas correntes de resposta : • os que aquiescem positivamente à integração ; • os que rejeitam, acerbamente, tal possibilidade.
No meio, uma corrente que somente admite a integração com modificações na conduta dos integrantes do movimento, com a exclusão de letras exclusivamente com apelos sexuais e incitação à violência.
A interpretação positiva indica a origem da música na periferia, notadamente nas comunidades cariocas e sua projeção para o Brasil e o mundo.
Nessa perspectiva, surge o reconhecimento como cultura popular comparada ao samba, axé e outros ritmos também nascidos no Brasil.
Também ocorre comparação com o movimento hip-hop americano.
Já a análise negativa prioriza a vinculação do ritmo com traficantes, prostitutas, vilipêndio às mulheres, falta de controle quanto ao som, local dos eventos, transações sexuais em público, uso frequente de drogas, violência exacerbada.
Faltam, para essa corrente, elementos básicos da cultura, tratando -se de música sem autonomia como movimento de mudança ou aspecto próprio.
A corrente intermediária reconhece o vínculo originário com as comunidades dominadas por quadrilhas de traficantes, mas admite o seu reconhecimento de acordo com a mudança dos atores que transitam na modalidade, como músicos, intérpretes, para adequar as letras das canções expurgando os aspectos considerados nocivos como o apelo sexual, as indicações positivas para a violência e apologia às drogas.
Um aspecto bastante criticado nas músicas dessa modalidade foi a exploração da mulher, entendida como um objeto de satisfação masculina.
É relevante tal constatação, pois a questão pertinente à Malala reflete a luta pela valorização da mulher.
Os concluintes compreenderam a ideia que foi problematizada e escolheram uma estratégia adequada para solucioná -la.
O caráter dinâmico da cultura foi comumente expressado, não através do enfrentamento abordado no texto, mas mediante a indicação ( de forma rasa e genérica ) das dificuldades e consequente aceitação social de manifestações culturais, antes marginalizadas e que hoje integram o acervo cultural nacional, tais como : samba, grafite, capoeira e forró.
Quanto às reflexões de Laraia, a grande maioria reconheceu o funk como manifestação cultural, porém, o identificou mais na sua expressão musical, no balanço de seu ritmo, tecendo contundentes críticas, e mesmo forte discordância quanto ao conteúdo e mensagem das letras quando remetem à depreciação da imagem feminina ou induzem à apologia ao crime. de seu ritmo, tecendo contundentes críticas, e mesmo forte discordância quanto ao conteúdo e mensagem das letras quando remetem à depreciação da imagem feminina ou induzem à apologia ao crime.
De resto, os textos abordaram bem mais o item dois do padrão de Respostas.
Uma minoria enfrentou os temas mudança e dinamismo cultural.
As respostas da corrente positiva indicam que o preconceito sempre acompanhou as manifestações populares como o samba, na sua origem.
Admitem que diante da miscigenação surgem movimentos culturais de várias naturezas, o que inclui o funk.
Também indicam a dificuldade de lazer dos membros das comunidades o que faz surgirem eventos para preencher essa necessidade.
Nesse âmbito, apontam a música como alternativa de ascensão social e integração aos bens de consumo.
As respostas da corrente negativa vinculam a música ao tráfico de drogas cujos elementos comandam as comunidades, propiciando que os bailes sejam mecanismos de venda e consumo dessa mercadoria ilícita.
Na mesma toada, indicam a vinculação com a exploração sexual, principalmente das mulheres maiores e menores de idade.
Indicam aspectos imorais nas letras.
Apresentam os locais dos bailes como inadequados e provocadores da desordem pública.
Recusam que músicas com discriminação às mulheres sejam consideradas como cultura.
A corrente intermediária indica a possibilidade de adoção do funk como cultura desde que ocorram modificações nas letras e no comportamento dos atores ( músicos, letristas e público ), retirando os aspectos negativos já apontados.
Fala -se em “ funk do bem ”.
A maior parte das respostas apresentadas indica que os concluintes têm conhecimento dos conflitos sociais e, no caso da questão do funk, conseguem identificar a origem popular da música.
A partir daí é que ocorre a cisão entre os que defendem o movimento como cultura, aduzindo a ascensão social dos artistas envolvidos dos que vinculam a música com atos ilícitos, rejeitando sua inserção no âmbito cultural.
Em algumas provas, há dificuldade de concatenação lógica de raciocínio, mas a maior parte consegue desenvolver o texto de forma razoável.
Essa questão permite visualizar a divisão social com os pensamentos liberais e conservadores, sendo bastante claros nas respostas dos concluintes. 3.3.1.5 Análise de Língua Portuguesa das Questões Discursivas do Componente de Formação Geral Os dados de Relações Internacionais, obtidos a partir das respostas às questões discursivas do Componente de Formação Geral, no que tange à Língua Portuguesa, encontram -se na Tabela 3.16 e no Gráfico 3.9.
Nesse aspecto, os alunos, de todo o Brasil, tiveram Média 64,3.
A maior Média com respeito à Língua Portuguesa foi obtida na região Norte ( 66,0 ), e a menor, na região Sul ( 62,4 ).
Quanto à variabilidade das notas, o Desvio padrão de todo o Brasil foi 22,6.
O menor Desvio padrão foi obtido na região Norte ( 17,5 ) e o maior Desvio padrão foi obtido na região Sul ( 24,7 ).
A Mediana das notas de Língua Portuguesa foi 69,5, a mesma obtida nas regiões Nordeste e Sudeste, enquanto nas demais regiões do Brasil a Mediana foi : 69,0 nas regiões Norte e Centro-Oeste ; e 68,5 na região Sul.
A nota Máxima para todo o Brasil foi de 100,0, com, pelo menos, um aluno tirando essa nota na região Sudeste.
Além disso, a nota Mínima foi zero em todas as regiões do país, sem exceção.
O Gráfico 3.9 mostra a distribuição das notas de Língua Portuguesa do Componente de Formação Geral.
Observa -se que a maior frequência corresponde aos alunos que obtiveram nota no intervalo ( 70 ; 80 ]. 3.3.1.6 Comentários sobre a correção das respostas de Formação Geral com respeito à Língua Portuguesa Ao encaminhar as questões 1 e 2 na direção da produção de um texto dissertativo, espera -se que o participante utilize seus conhecimentos sobre o assunto e estruture seus textos de acordo com as exigências do registro formal próprio dessa situação comunicativa.
Essa configuração determina exigências quanto : à adequação da seleção vocabular, ao desenvolvimento do conteúdo, à estruturação sintática dos períodos, à organização lógica das ideias, à utilização de procedimentos de encadeamento textual e referenciação, à obediência às exigências morfossintáticas próprias da modalidade escrita da norma-padrão, ao respeito às regras ortográficas e às regras de acentuação gráfica.
O padrão de resposta utilizado na avaliação das questões 1 e 2 considerou os aspectos relevantes ao bom desempenho linguístico como competências distintas, de modo a permitir um mapeamento detalhado do domínio dos recursos disponíveis na Língua Portuguesa para a comunicação escrita formal.
Com base nesse objetivo, foram avaliados os seguintes aspectos : a ) Estruturação textual condizente com o gênero solicitado ( texto dissertativo ) e o modo de organização textual expositivo adequado ao gênero – essa competência envolve :  distribuição do conteúdo do texto em parágrafos, de modo a garantir a sua organização temática ;  estruturação sintática condizente com o padrão da modalidade escrita formal da língua portuguesa de modo a garantir a clareza necessária ;  utilização de operadores discursivos que contribuam para a progressão temática do texto, estabelecendo relações lógicas entre as ideias apresentadas, tanto do ponto de vista intrafrasal, como do interfrasal ;  utilização de procedimentos de referenciação lexical e pronominal que permitam a retomada de referentes textuais ;  utilização de sinais de pontuação que contribuam para a organização lógica da frase e do texto.
Espera -se, portanto, que o participante recorra a procedimentos linguístico-discursivos para organizar seu texto, permitindo o encadeamento lógico entre suas partes de forma a garantir a progressão e a coerência textual.
Isso significa que os seguintes procedimentos foram penalizados, de acordo com o padrão de resposta proposto :  elaboração de frases fragmentadas que comprometam a estrutura lógico-gramatical do texto ;  sequência justaposta de ideias sem encaixamentos sintáticos, reproduzindo hábitos da oralidade ;  elaboração de frase com apenas oração subordinada, sem oração principal ;  emprego equivocado do conector ( preposição, conjunção, pronome relativo, alguns advérbios e locuções adverbiais ) que não expresse a relação lógica adequada entre dois trechos do texto e prejudique a compreensão da mensagem ;  repetição ou substituição inadequada de palavras sem se valer dos recursos oferecidos pela língua ( pronome, advérbio, artigo, sinônimo ) ;  utilização inadequada dos sinais de pontuação comprometendo a clareza textual. b ) Respeito às convenções ortográficas da norma-padrão da Língua Portuguesa – essa competência envolve o domínio das regras de acentuação gráfica e da grafia padrão das palavras ( com ausência de abreviaturas próprias da linguagem da internet ), de acordo com as convenções estabelecidas pela legislação em vigor e consubstanciadas no Vocabulário Ortográfico da Língua Portuguesa, editado pela Academia Brasileira de Letras ( com aceitação da legislação anterior, no caso das regras relativas ao uso do hífen e da acentuação gráfica, já que o exame foi realizado no mês de novembro de 2015 ).
A competência relativa ao respeito às convenções ortográficas envolve :  grafia correta das palavras ;  respeito às regras de acentuação gráfica ;  emprego de maiúsculas em início de frase, em nomes próprios de pessoas, lugares ou instituições ;  ausência de abreviações como p/, vc, tb, pq, \textbf{tá}, né, usadas muitas vezes em escrita informal e na internet ;  obediência às regras de separação de sílabas no final da linha. c ) Domínio dos diferentes aspectos morfossintáticos próprios da modalidade escrita formal da norma-padrão da Língua Portuguesa – os itens relacionados a essa competência são : a concordância nominal, a concordância verbal, a regência nominal, a regência verbal, a flexão nominal, a flexão verbal, a correlação entre os tempos verbais, a colocação pronominal.
O domínio dessa competência envolve :  flexão do verbo para estabelecer concordância de número com o sujeito da frase ;  flexão do artigo, do adjetivo e do pronome para concordar em número e em gênero com o substantivo a que se referem ;  respeito à regência nominal e à verbal, utilizando a preposição adequada depois de um substantivo, de um verbo ou de um adjetivo ; emprego da preposição antecedendo o pronome relativo e emprego adequadamente os pronomes relativos “ cujo(a) ” e “ onde ” ;  emprego adequado do acento grave indicador da crase entre uma preposição e um artigo ( a+a ) ;  atendimento às regras de colocação pronominal ( próclise e ênclise ), distintas dos hábitos da oralidade ou da escrita informal. d ) Seleção vocabular adequada à modalidade escrita formal da Língua Portuguesa, exigida pela situação comunicativa – essa competência envolve a precisão na utilização do vocabulário relacionado à temática solicitada pela questão ; a ausência de marcas da oralidade, como termos de sentido muito genérico ( “ coisa ”, “ negócio ”, “ você ” ) e termos de registros mais informais ( como gírias, jargões, frases feitas, ditados populares, termos regionais ).
Assim, espera -se que o participante respeite a adequação vocabular não empregando gírias ou expressões coloquiais, evite repetição desnecessária de palavras e utilize um vocabulário mais formal, como solicitado por um texto dissertativo.
A escolha dessas competências para subsidiar o processo de avaliação apoia -se na concepção de que, no desempenho dos graduandos, a modalidade escrita tem apresentado uma intensa simplificação, originada no padrão da modalidade oral da Língua Portuguesa.
No caso do texto de base dissertativa, inscrito em um registro formal, a distância entre as duas modalidades é ainda maior, o que provoca situações de hipercorreção ( desvios provocados pela utilização inadequada de uma regra da norma-padrão ) e de truncamentos sintáticos ( estruturas frasais incompreensíveis devido à complexidade sintática própria da modalidade escrita ).
Observam -se, então, os seguintes aspectos que marcam essa distinção entre as duas modalidades, devido à excessiva simplificação da modalidade falada : a ) redução drástica de estruturas subordinadas, compensada pelo aumento na frequência de estruturas coordenadas e absolutas, por um lado, ou pela elaboração de estruturas truncadas pelo excesso de ideias sem a devida conexão subordinativa ; b ) redução no uso de operadores argumentativas para expressar relações lógicas essenciais à construção do texto, substituídas pela exigência de inferência por parte do interlocutor para suprir a sua ausência ; c ) redução do uso do subjuntivo, ao lado da ampliação do uso do indicativo combinado a estruturas frasais coordenadas ou absolutas ; d ) empobrecimento do processo de referenciação, com a repetição de pronomes ou nomes ; e ) simplificação extrema da marcação da categoria tempo na morfologia verbal ; f ) falta de domínio de vocabulário mais abstrato e de maior complexidade, essencial ao desenvolvimento do processo dissertativo ; g ) redução drástica no emprego da acentuação gráfica, processo intensificado pela intensa utilização das redes sociais.
Os aspectos macroestruturais da elaboração do texto não foram avaliados neste processo porque dizem respeito à avaliação do conteúdo.
Portanto, a banca de Formação Geral, composta por profissionais de diferentes áreas do conhecimento, se encarregou da avaliação do atendimento ao solicitado no enunciado das questões do ponto de vista do desenvolvimento do conteúdo.
A avaliação do desempenho linguístico considerou, portanto, três grandes grupos de competências, segundo os aspectos explicitados anteriormente : a ) domínio das convenções ortográficas : grafia de vogais e consoantes, uso de maiúsculas e minúsculas, emprego do hífen e acentuação gráfica ; b ) domínio dos procedimentos de estruturação textual do ponto de vista microestrutural : organização interna dos períodos, emprego de conectores para a articulação lógica entre os períodos e entre os parágrafos, emprego de marcas de referenciação lexical e pronominal ; utilização dos sinais de pontuação que contribuem para a organização lógica da frase. c ) domínio das regras de caráter morfossintático estabelecidas como modelares do ponto de vista da modalidade escrita formal da norma-padrão da Língua Portuguesa : concordância nominal e verbal, regência nominal e verbal, colocação pronominal, flexão nominal e verbal, correlação entre tempos e modos verbais, ausência de marcas de oralidade.
A seleção vocabular adequada à modalidade escrita formal da Língua Portuguesa foi incorporada à essa última competência, tendo em vista a intersecção entre as duas do ponto de vista das exigências do registro formal da modalidade escrita da norma-padrão.
Apresenta -se, a seguir, o padrão de resposta aprovado pelo Inep e já utilizado na avaliação do desempenho linguístico das questões 1 e 2 de Formação Geral no Enade/2013 e no Enade/2014.
Às competências, reunidas nos três grupos descritos anteriormente, foram atribuídos os seguintes pesos relativos : aspectos ortográficos ( 20 \% ) ; aspectos textuais ( 40 \% ), aspectos morfossintáticos e vocabulares ( 40 \% ).
Com base na avaliação do processo desenvolvido nos dois anos anteriores, utilizaram -se quatro níveis para a avaliação do desempenho linguístico nas três competências.
O nível zero ficou reservado para casos especiais, como fuga total ao tema, palavras soltas, produção de uma frase incompleta, entre outros.
A avaliação das respostas elaboradas pelos participantes revelou um resultado coerente na comparação entre as duas questões.
Entretanto, em relação aos resultados do Enade/2014, observou -se uma melhora do domínio de todas as competências, talvez devido ao perfil dos cursos analisados, como Direito e Jornalismo, que pressupõem a utilização profissional da linguagem.
Observou -se uma diferença no desempenho dos participantes em relação às duas questões discursivas, provavelmente em função dos temas solicitados : - a Questão 1 solicitava que o participante desenvolvesse um texto dissertativo sobre o significado da premiação de Malala, abordando dois aspectos : o direito das jovens à educação formal e as relações de poder entre homens e mulheres no mundo.
Esse encaminhamento favoreceu o desenvolvimento mais articulado do texto, o que se refletiu em melhor desempenho dos aspectos vocabulares e textuais. - a Questão 2 solicitava que o participante desenvolvesse um texto dissertativo para discutir a questão do reconhecimento do funk como legítima manifestação artística e cultural do povo brasileiro.
Por ser um tema muito polêmico, o tema do preconceito sócio-cultural-econômico favoreceu a construção de textos mais críticos, dividindo o conjunto de participantes em dois grupos : os que defendiam esse tipo de manifestação artística como legítima e os que a criticavam de maneira veemente, chegando, inclusive, a utilizar muitos termos de baixo calão.
Essa última atitude teve como consequência a elaboração de textos mal desenvolvidos do ponto de vista sintático-discursivo, com muitos truncamentos.
Os resultados revelaram essa diferença.
Ao analisar os dados de \textbf{cada} questão, constata -se que a Questão 1 teve maior número de questões de médio ou bom desempenho, em relação à Questão 2.
Por outro lado, a Questão 2 apresentou um índice maior de notas baixas em relação à Questão 1.
Os resultados da avaliação correspondem aos seguintes aspectos observados em \textbf{cada} competência : a ) aspectos ortográficos : o desempenho dos participantes revelou uma diferença muito grande nos dois aspectos analisados : baixo índice de desvios ortográficos e grande índice de desvios de acentuação.
Em vários casos, ocorre ausência completa de acentuação gráfica.
Os resultados revelam que a tendência dominante entre os universitários brasileiros é a eliminação da acentuação gráfica, talvez motivada pelos hábitos relacionados às redes sociais e pela ausência de esclarecimento dos meios de comunicação, das autoridades e das escolas sobre as decisões do Acordo Ortográfico da Língua Portuguesa de 1990.
Os casos mais sistemáticos de eliminação do acento indicador da sílaba tônica são :  palavras proparoxítonas ( “ proximos ”, “ politicos ” ) ;  palavras paroxítonas terminadas em ditongo crescente ( “ necessario ”, “ noticiarios ”, “ individuo ”, “ dependencia ”, “ varios ”, “ propria ” ) ;  palavras oxítonas ( “ ninguem ”, “ esta ”, “ ate ”, “ tambem ”, “ alvara ” ) ;  uso indevido do acento gráfico em palavras como “ genêro ”, “ melâncolia ”, por exemplo. exemplo.
Quanto ao domínio das convenções relativas à grafia das palavras, observam -se desvios como : I. hipercorreção pela escolha de “ e ” no lugar de “ i ” ou de “ i ” no lugar de “ e ”, por influência de hábitos da oralidade ( “ descriminação ” por “ discriminação ”, “ entervenção ” por “ intervenção ”, “ esteriotipada ” por “ estereotipada ” ) ; II. eliminação do “ r ” marcador do infinitivo verbal ( “ esta ” no lugar de “ \textbf{estar} ”, “ estuda ” no lugar de “ estudar ” ).
III. eliminação de uma sílaba na palavra “ educação ”, que foi grafada “ edução ” por um grande número de participantes.
IV. outros casos de desvios de grafia relacionados à variação diastrática podem ser observados em “ estrupo ”, “ subjulgadas ”, “ precoseito ”, “ soubre ”, “ vecendo ”, “ indesencia ”, “ apolojia ”, “ fulga ”.
V. desvios de segmentação podem ser atribuídos à confusão entre o morfema gramatical “ mos ” e o pronome oblíquo “ nos ”, em “ chegar mos ”, “ enfrentar -mos ”.
VI. inadequação no uso da maiúscula para destacar determinadas palavras-chave do texto, como “ Violência ”, “ Brasileiros ”.
Destaque -se, também, o grande número de participantes que grafam os textos inteiramente em caixa alta.
Vale observar, também, que, ao contrário do que se esperava, não apareceram abreviaturas próprias do “ internetês ”, relacionado ao uso de redes sociais e emails. b ) aspectos textuais : esta competência é a que se revela como a mais problemática entre os participantes, porque são muitos os problemas observados, desvios acumulados durante toda a formação do estudante : sequência justaposta de ideias sem encaixamentos sintáticos ; redução drástica de estruturas subordinadas, ao lado do aumento na frequência de estruturas coordenadas e absolutas ; redução no uso de conectores para expressar relações lógicas essenciais à construção do texto, substituídas pela exigência de inferência por parte do interlocutor para suprir a sua ausência ; emprego equivocado de operadores que não estabelecem relações lógicas coerentes entre ideias do texto ; emprego inadequado do pronome relativo ( com omissão da preposição ou a utilização de pronome inadequado, como “ onde ” ) ; repetição exaustiva de termos sem a utilização de procedimentos mais sofisticados de substituição ( hiperonímias, hiponímias, nominalizações, expressões metafóricas ) ; frases fragmentadas que comprometem a estrutura lógico-gramatical ; frases formadas apenas por oração subordinada, sem oração principal. fragmentadas que comprometem a estrutura lógico-gramatical ; frases formadas apenas por oração subordinada, sem oração principal.
Um importante aspecto a destacar é o baixíssimo desempenho de uma parte dos participantes em relação à estrutura formal do texto produzido, o que é extremamente preocupante ao se levar em conta que são graduandos em fase final de formação.
São frequentes os casos de desvios de estruturação frasal, com uso inadequado ou ausência de conectivos entre parágrafos e entre frases.
Em uma parte dos textos, falta um mínimo de textualidade e de domínio do registro padrão da língua.
Na verdade, observam -se relações linguísticas quase agramaticais, como as estabelecidas pela sequência de gerúndios sem o apoio de um ponto de partida para a organização das informações gramaticais e semânticas.
Observou -se que uma grande parte dos participantes não organizou as ideias em parágrafos, talvez devido ao pequeno número de linhas disponibilizadas para a resposta da questão ou, quem sabe, pela suposição de que não seria necessária essa divisão por não se tratar de um texto no modelo de uma redação dissertativo-argumentativa, como solicitado nos vestibulares.
Quanto à utilização dos mecanismos de referenciação, deve -se destacar a ocorrência, em uma boa parte dos textos, de repetições de palavras ou expressões sem a utilização de termos sinônimos ou pronomes, como seria adequado.
Outro aspecto relevante a ser destacado é a quase total ausência de operadores argumentativos, tanto intrafrasais como interfrasais, repercutindo uma tendência atual da mídia escrita.
Quanto à utilização dos sinais de pontuação, observou -se uma grande precariedade nos textos analisados.
É muito frequente a ocorrência de parágrafos sem marca interna de pontuação para separar os períodos.
Vale observar que não foi penalizada a ausência de vírgula para destacar locuções ou adjuntos adverbiais de pequena extensão deslocados de posição na frase, por ser um uso opcional.
São os seguintes os tipos de problemas encontrados : I. vírgula : utilização de vírgula para separar o sujeito e o predicado ; ocorrência de apenas uma das vírgulas para separar uma palavra, uma expressão ou uma oração encaixada ; uso de vírgula no lugar do ponto para separar ideias que constituem períodos distintos ; II. ponto e vírgula : utilização do ponto e vírgula no lugar de vírgula ; III. ponto final : ausência de ponto final para separar períodos. c ) aspectos morfossintáticos e vocabulares : os resultados são muito transparentes em relação aos aspectos mais problemáticos no desempenho dos participantes.
O desvio mais frequente, em relação à regência, é a falta do sinal indicativo da crase – isso revela que o usuário não tem consciência de que, sob a forma do termo “ a ”, existe a presença de uma preposição “ a ”, exigida pela regência do termo anterior, e de um artigo definido.
Embora em outros exames, como o Enem, a falta de crase seja penalizada em acentuação, nesta avaliação esse desvio foi considerado no âmbito dos aspectos morfossintáticos.
Outro problema relacionado à regência verbal e à nominal, encontrado frequentemente nas questões, foi a ausência de preposição antes de pronome relativo, processo generalizado na modalidade oral da língua, em situações de registro informal.
Apesar da possibilidade de que essa alteração de regência se generalize no padrão escrito da Língua Portuguesa, como já \textbf{está} ocorrendo até em textos jornalísticos, o não emprego da preposição foi penalizado neste processo de avaliação.
Outro desvio muito frequente diz respeito aos processos de concordância verbal e de concordância nominal.
Quanto à concordância de número, observou -se ausência de marca com sujeito posposto ou uso indevido ( uso inadequado da marca de plural comandado pelo núcleo plural da locução adjetiva, apesar de o substantivo que funciona como núcleo do sintagma nominal \textbf{estar} no singular ).
Uma ocorrência generalizada foi a ausência de acento circunflexo na forma plural do presente do indicativo do verbo “ ter ”, que foi considerada como um desvio na concordância verbal e não na acentuação gráfica.
Quanto à concordância de gênero, vários casos foram observados, normalmente no âmbito de sintagmas nominais longos, em que o adjetivo \textbf{está} afastado do substantivo.
Quanto à questão da colocação pronominal, foram poucos os desvios observados.
Apesar de serem aspectos relacionados à oralidade, concluiu -se que, no registro escrito formal, a maioria dos participantes já incorporou regras como a não introdução da frase por um pronome oblíquo e a próclise na presença de um termo atrator.
Não se adotou, entretanto, o padrão excessivamente formal descrito pelas gramáticas normativas em relação à posição do pronome oblíquo em locuções verbais, já que esse uso \textbf{está} muito distante da prática cotidiana, até em textos mais formais.
Quanto aos aspectos vocabulares, alguns tipos de inadequação foram observados, principalmente na Questão 2, relacionada ao funk : expressões da oralidade, seleção vocabular incompatível com o contexto, gerando situações de falta de inteligibilidade ; falta de domínio de vocabulário mais abstrato e de maior complexidade, essencial ao desenvolvimento do texto de base dissertativa.
Ainda em função da temática desenvolvida na Questão 2, foi significativa a presença de palavras de baixa calão, por retratar um universo cultural específico.
Os textos que continham esse tipo de vocabulário foram considerados inadequados em função do registro formal exigido e foram avaliados com grau zero.
Várias marcas de oralidade foram identificadas, embora não em alta frequência : o uso do pronome relativo “ onde ” como relativo universal, falta de artigo definido antes de substantivo, repetição de palavras por falta de vocabulário, reduções como “ \textbf{tá} ” e “ pra ”, expressões informais. 3.3.2 Componente de Conhecimento Específico Na parte da prova relativa às questões discursivas do Componente de Conhecimento Específico ( Tabela 3.17 ), observa -se que a Média foi bem mais baixa do que para as questões discursivas do Componente de Formação Geral.
Enquanto no Componente de Formação Geral a Média para estudantes de Relações Internacionais de todo o Brasil foi 64,6, na parte de Conhecimento Específico a Média foi 34,0.
A maior Média deste componente foi obtida pelos estudantes da região Sudeste ( 35,1 ), e a menor, pelos da região Norte ( 26,0 ).
Quanto à variabilidade das notas, o Desvio padrão de todo o Brasil foi 22,8.
O maior Desvio padrão foi encontrado nas regiões Sul e Centro-Oeste ( 23,4 ), e o menor, na região Norte ( 19,7 ).
A maior nota Máxima foi obtida na região Sudeste ( 96,7 ), enquanto a menor nota Máxima foi encontrada na região Nordeste ( 80,0 ).
Além disso, a nota Mínima ( 0,0 ) foi obtida por pelo menos um aluno em todas as regiões do Brasil, sem exceção.
A Mediana do Brasil como um todo foi 35,0, a mesma obtida na região Sul.
Nas demais regiões a Mediana foi : 23,3 na região Norte ; 31,7 na região nordeste ; 36,7 na região Sudeste ; e 30,0 na região Centro-Oeste.
O Gráfico 3.10 representa a distribuição das notas nas questões discursivas no Componente de Conhecimento Específico.
A moda desta distribuição ocorre no primeiro intervalo, [ 0 ; 10 ].
Nota -se que a distribuição é decrescente a partir do intervalo ( 30 ; 40 ].
A análise de \textbf{cada} uma destas questões será feita a seguir. 3.3.2.1 Análise da Questão Discursiva 3 do Componente de Conhecimento Específico Na questão 3, cujos resultados aferidos encontram -se descritos na Tabela 3.18, a Média dos estudantes de todo o Brasil foi 34,0.
A menor Média nessa questão foi obtida pelos alunos da região Norte ( 30,1 ), enquanto a maior Média foi obtida na região Nordeste ( 35,4 ).
Quanto à variabilidade das notas, o Desvio padrão de todo o Brasil foi 31,3.
O maior Desvio padrão foi obtido na região Centro-Oeste ( 32,5 ), enquanto o menor foi obtido na região Norte ( 29,8 ).
As medianas ( 30,0 ), as notas mínimas ( 0,0 ) e as notas máximas ( 100,0 ) foram as mesmas em todas as regiões do Brasil, exceto a Mediana na região Nordeste ( 35,0 ).
O Gráfico 3.11 mostra a distribuição das notas na questão discursiva 3, do Componente de Conhecimento Específico da área de Relações Internacionais.
Essa distribuição é multimodal com moda principal no intervalo ( 60 ; 70 ]. 3.3.2.2 Comentários sobre as respostas à Questão Discursiva 3 O objetivo principal da questão foi o de avaliar habilidades de compreensão da Língua Inglesa, fundamental para uma formação em Relações Internacionais, considerando -se o caráter internacional da Área que deve formar o profissional capaz de ter acesso à literatura adequada, a fontes de informação ou a necessidade da inserção e interação internacional do profissional.
Embora não seja cobrada a fluência escrita do inglês nos cursos da Área pressupõe -se uma mínima capacidade de o graduado interpretar textos simples.
Dessa forma, não foi exigida resposta em inglês, foi apenas esperada a compreensão de um texto simples, sobre um conteúdo fácil, que deveria se refletir na resposta dissertativa em português.
Apesar da facilidade da questão, houve alto índice de notas zero.
Como já mencionado, a resposta ao item ‘ a ’ \textbf{estava} contida, mesmo que parcialmente, no próprio texto do enunciado.
Assim, pudemos concluir que o baixo desempenho é fortemente marcado pela falta de compreensão da língua inglesa escrita.
Em alguns casos, apesar de demonstrarem compreensão do texto, os alunos não apresentaram objetividade nas suas respostas, no sentido de demonstrar o conhecimento pretendido.
Discorreram genericamente sobre a presença da Missão das Nações Unidas para Estabilização do Haiti - Minustah, por exemplo -, sem atender objetivamente ao que foi solicitado, de forma clara, e \textbf{estava} previsto pelo padrão de resposta.
Em geral, os alunos adotaram linguagem adequada da Língua Portuguesa em suas respostas, mas como ocorreu em outras questões, uma minoria de alunos demonstrou sérias dificuldades com o português e sua capacidade argumentativa.
O desempenho geral foi baixo, mas vale destacar a existência de um gap, uma disparidade entre excelentes respostas e respostas muito fracas.
Nessa questão, embora haja uma régua de avaliação adequada e capaz de capturar acertos parciais, pode -se conferir uma polarização entre notas acima de 6,0 e as abaixo de 4,0 ( incluindo o zero ).
Frequentemente as respostas que receberam notas muito baixas evidenciavam dificuldade de os alunos desenvolverem argumentos.
Por exemplo, ao buscarem responder ao comando, muitos resumiam a “ motivação doméstica ” na “ busca de uma vida melhor ”.
As respostas foram frequentemente fracas ou sem foco em relação ao que foi solicitado.
Muitas fizeram referências genéricas ao Brasil como um país receptivo e “ amigo ”, que protege pobres e afrodescendentes, sempre sem fundamentação adequada ou conexão com o conteúdo esperado.
Mas, considerando -se o objetivo de 3.3.2.3 Análise da Questão Discursiva 4 do Componente de Conhecimento Específico A Tabela 3.19 contém as informações relativas à questão 4 do conjunto de questões do Componente de Conhecimento Específico.
O desempenho dos estudantes de todo o Brasil nesta questão foi mais baixo dentre as três questões discursivas desse componente.
A Média geral do Brasil foi 23,7, sendo a menor Média registrada na região Centro-Oeste ( 18,8 ), e a maior, na região Sul ( 25,3 ).
As medianas ( 20,0 ), as notas mínimas ( 0,0 ) e as notas máximas ( 100,0 ) foram as mesmas em todas as regiões do Brasil, sem exceção. aferir compreensão da Língua Estrangeira, foram aproveitadas, ao máximo, respostas que demonstraram capacidade de compreensão do texto em inglês.
Respostas fracas frequentemente ignoravam a primeira parte do comando e mencionavam “ terremotos ou desastres ambientais ” como causas genéricas para a imigração de haitianos para o Brasil – sem considerarem o “ turning point ” referido na pergunta.
Algumas respostas muito fracas mencionaram furacões e guerras no Haiti, e outras se referiam ao Haiti como país africano e sua língua como o português.
Pode -se dizer que houve pouca variedade, tanto nos erros quanto nos acertos e não trouxeram nenhuma dificuldade para a aplicação do padrão de resposta.
Certamente, a conclusão geral foi a de que o conhecimento da Língua Inglesa apresentado pelos concluintes é baixo.
Apesar de não ser cobrada a proficiência na Língua Inglesa ou a habilidade da escrita, uma mínima compreensão de leitura e interpretação de texto é fundamental para a Área, e a questão revela uma necessidade de maior atenção a esse ponto.
Resta ressaltar que algumas poucas respostas foram elaboradas em inglês.
Embora o comando da questão solicitasse a resposta em português.
Apesar disso, as respostas apresentadas em inglês bem escrito e de conteúdo correto foram consideradas.
Houve raros casos de respostas anuladas por uso de linguagem inadequada e alguns protestos. 3.3.2.3 Análise da Questão Discursiva 4 do Componente de Conhecimento Específico A Tabela 3.19 contém as informações relativas à questão 4 do conjunto de questões do Componente de Conhecimento Específico.
O desempenho dos estudantes de todo o Brasil nesta questão foi mais baixo dentre as três questões discursivas desse componente.
A Média geral do Brasil foi 23,7, sendo a menor Média registrada na região Centro-Oeste ( 18,8 ), e a maior, na região Sul ( 25,3 ).
As medianas ( 20,0 ), as notas mínimas ( 0,0 ) e as notas máximas ( 100,0 ) foram as mesmas em todas as regiões do Brasil, sem exceção.
O Gráfico 3.12 representa a distribuição de notas da questão discursiva 4, do Componente de Conhecimento Específico.
Essa distribuição é multimodal com moda principal no intervalo ( 10 ; 20 ]. 3.3.2.4 Comentários sobre as respostas à Questão Discursiva 4 A questão visava, tal como ficou refletido no padrão, a respostas simples e diretas que deveriam apresentar informações objetivas, sem exigência de texto argumentativo.
Procurou -se, mesmo em casos problemáticos, evitar prejuízo ao aluno, buscando verificar a presença dos aspectos elencados no padrão de resposta no texto redigido pelos alunos.
O conteúdo da questão pode ser considerado fácil por ser abordado em disciplinas obrigatórias, presentes em qualquer curso de Relações Internacionais.
Apesar disso, a objetividade e facilidade não se refletiu na atribuição de notas.
Houve um grande número de respostas que fugiam à questão e apelavam para argumentações de senso comum e ideológico.
Uma observação interessante que refletiu a má performance, em geral, dos estudantes, foi que o espaço designado para a resposta não foi devidamente aproveitado.
Houve um número razoável de respostas curtas, quando objetivas, mas não necessariamente corretas ou frases abrangentes de caráter ideológico ou com expressa opinião do estudante, sem relação com os conceitos e os conteúdos de Política Externa Brasileira ( PEB ) esperados pelo padrão de resposta.
A régua de pontuação foi adequada para a correção, mas o descompasso entre o padrão de resposta e as respostas encontradas, na maioria dos casos, se deu por quatro razões : i ) confusão sobre propósitos da PEB e da política doméstica ; ii ) confusão entre as políticas dos dois presidentes citados ; iii ) falta de base em PEB ; iv ) conhecimento insuficiente de comércio exterior ou história recente do Brasil.
Em geral, os alunos adotaram a forma adequada da Língua Portuguesa em suas respostas, mas como ocorreu em outras questões há aqueles, apesar de ser uma minoria, que demonstraram sérias dificuldades no uso da norma culta da língua e em sua capacidade argumentativa.
O desempenho geral foi inesperadamente fraco, considerando -se que o conteúdo pertence à disciplina obrigatória sobre Política Externa Brasileira nos cursos de Relações Internacionais.
Isso evidenciou uma baixa capacidade de elaborar argumentos e atribuir conceitos específicos e técnicos referentes ao conteúdo da questão ; consequentemente, a correção revelou deficiências na formação específica dos alunos como também insuficiência no domínio da Língua Portuguesa.
Dentre as respostas fracas, as mais comuns foram aquelas mal elaboradas ou com argumentos de senso comum e/ou ideológico ou que, em alguma parte, mencionava a intensidade das relações Sul-Sul na política externa do período do governo de Luís Inácio Lula da Silva.
Dentre as respostas medianas destacam -se aquelas nas quais os concluintes conseguiram responder de forma completa ao item ‘ b ’, mas não totalmente ao item ‘ a ’.
Isso revelou claramente a falta de conhecimento sobre a agenda comercial externa brasileira.
O item ‘ b ’ permitia o apelo a informações mais amplamente conhecidas sobre a atuação externa dos governos, ou seja, trata de conteúdo menos técnico e específico que se espera de um profissional da área de Relações Internacionais.
Sobre a formação da Área destaca -se uma surpreendente falta de domínio e de informação sobre a política externa contemporânea.
Como já comentado, muitas respostas se baseavam no senso comum e/ou ideológico, talvez como tentativa de redigir alguma resposta geral quando não havia domínio sobre o conteúdo.
Em resumo, o desconhecimento do conteúdo e a frequente confusão entre o conteúdo de política externa, da disciplina de PEB, com opiniões partidárias e ideológicas, revelaram a falta de compreensão sobre o significado de uma prova que deveria mensurar a aprendizagem específica do ensino superior na Área.
Além disso, o uso de um linguajar não apropriado e a incapacidade de desenvolvimento de argumentos demonstraram, na verdade, falta de domínio de debates, fatos, história e conceitos considerados fáceis para a Área.
Houve raros casos de anulação da resposta por uso de linguagem inadequada e alguns protestos.
Dentre os protestos mais comuns, \textbf{estão} aqueles a respeito do sentimento contrariado em realizarem as provas do Enade/2015. 3.3.2.5 Análise da Questão Discursiva 5 do Componente de Conhecimento Específico A Tabela 3.20 contém as informações relativas à questão 5 do conjunto do Componente de Conhecimento Específico.
O desempenho dos estudantes nessa questão foi superior ao das questões 3 e 4.
A nota Média dos estudantes de todo o Brasil foi 44,4.
A maior Média foi registrada na região Sudeste ( 46,7 ), enquanto a menor Média foi registrada na região Norte ( 27,9 ).
Quanto à variabilidade das notas, o Desvio padrão dos alunos do Brasil, como um todo, foi 34,4.
Enquanto o maior desvio foi encontrado na região Centro-Oeste ( 35,8 ), o menor foi encontrado na região Norte ( 31,1 ).
A Mediana do Brasil como um todo foi 55,0, a mesma obtida na região Sudeste, enquanto nas demais regiões foi : 25,0 na região Norte ; 40,0 nas regiões Nordeste e Centro-Oeste ; e 50,0 na região Sul.
As notas mínimas ( 0,0 ) e as notas máximas ( 100,0 ) foram as mesmas em todas as regiões do Brasil, sem exceção.
O Gráfico 3.13 representa a distribuição das notas da questão discursiva 5 do Componente de Conhecimento Específico.
Destaca -se, o grande número de estudantes que deixaram esta questão em branco, correspondendo à moda da distribuição. 3.3.2.6 Comentários sobre as respostas à Questão Discursiva 5 O enunciado solicitava que o aluno conhecesse três princípios do Tratado de Não Proliferação Nucelar - TNP, porém poucos conseguiram citar os três princípios.
Geralmente, citavam a própria “ não proliferação ” e o desenvolvimento de tecnologia nuclear para fins pacíficos, citado no texto introdutório que compunha o enunciado.
Poucos citaram o desarmamento.
Outras respostas muito comuns e que não \textbf{estavam} entre os princípios do padrão de resposta foram “ manutenção da paz ”, “ segurança ” e “ preservação da humanidade ”.
Também foi possível notar que muitos estudantes não conseguiam conectar os argumentos favoráveis e desfavoráveis do TNP com a discussão sobre o Irã.
Nesse ponto, muitas vezes, a redação do texto prejudicava o conteúdo, que acabava por ficar muito confuso.
Tal situação levou a um maior esforço no sentido de avaliar o significado do que foi redigido e não a forma.
Muitas vezes, a linguagem adotada pelos respondentes também não foi adequada, no uso de termos específicos da Área, o que demonstra falta de leitura, incapacidade argumentativa e precariedade na redação.
Foi possível notar que alguns tiveram dificuldade de se ater ao enunciado e responder conforme o que foi solicitado pelo comando.
A questão pedia dois argumentos favoráveis e dois argumentos desfavoráveis, porém muitos apresentavam apenas um argumento favorável e, depois, repetiam o mesmo com outras palavras como se fosse um segundo argumento.
Também foi notado que alguns estudantes não apontavam os argumentos a serem avaliados como “ favorável ” ou “ desfavorável ” ao regime do TNP, deixando assim a desejar em relação ao esperado, segundo o padrão de resposta.
Outro ponto que merece comentário se refere ao comando do item ‘ a ’, que não delimita a quantidade de princípios do TNP que deveriam ser apresentados pelos respondentes.
Se por um lado, o comando foi redigido no plural – “ Aponte os princípios ” –, sem discriminação sobre quantos princípios se esperava que o aluno citasse, por outro, esta estrutura de redação poderia levar a crer que o estudante deveria apontar todos os princípios.
Para superar esse problema, em conjunto com a Comissão Assessora decidiu -se que dois princípios seriam suficientes para o aferimento da pontuação total, e a correção confirmou o acerto dessa decisão.
Os alunos com respostas fracas ou muito fracas foram minoria.
Em geral, esses não conheciam o conteúdo e apenas reproduziam o enunciado, de forma precária, apontando elementos ali presentes como argumento favorável e desfavorável.
As respostas medianas predominaram.
Nesses casos os alunos discorreram sobre o tema, demonstraram algum conhecimento, mas não dominavam plenamente o assunto, logo, omitiam parte da resposta esperada.
As respostas boas, como as fracas, representaram uma minoria.
No entanto, vale destacar que algumas respostas apresentavam qualidade muito boa, com textos bem desenvolvidos e argumentos adequados, tendo essas alcançado a pontuação máxima na correção.
Um número irrisório de concluintes declarou que não houve tempo para completar a prova.
No entanto, considerou -se que o espaço disponibilizado foi suficiente para o desenvolvimento da questão, pois alguns concluintes conseguiram ser sucintos e profundos nas respostas.
Considerando -se o padrão de resposta a ser seguido, foi constatado que a maior parte das respostas foram genéricas.
É possível que o fato de ser uma temática recorrente na imprensa tenha favorecido uma confusão, por parcela dos respondentes do ENADE, entre a discussão técnica da área e o modo como a polêmica é apresentada nos meios de comunicação. 3.3.3 Considerações Finais Em geral, as três questões específicas da Área de Relações Internacionais foram fáceis.
As questões 4 e 5, no entanto, tal como elaboradas e, conferindo -se o que se esperava do estudante pelo padrão de resposta, eram, de fato, questões objetivas.
Essa objetividade tornou difícil a adoção de uma régua de pontuação muito graduada, que refletisse aspectos mais detalhados do conteúdo da resposta.
Problemas com a Língua Portuguesa foram generalizados, não tanto quanto à ortografia, mas fundamentalmente foram encontradas grandes dificuldades de desenvolvimento de um texto argumentativo, bem como a obediência aos comandos das questões.
Muitos e generalizados foram os casos de respostas que escapavam aos objetivos das questões, não demonstrando preocupação em responder ao que foi solicitado.
Na questão 3, mesmo sem compreensão do inglês ou do conteúdo a ser conferido, um relevante número de respostas fracas apresentou uma visão ideológica sobre o Governo Lula como sendo um governo que distribuiu benefícios sociais aos pobres, negros e imigrantes.
Tema que escapa completamente ao que foi \textbf{perguntado}.
Na questão 4, sobre PEB, o posicionamento ideológico a favor ou contra dois governos que se opuseram por tantos anos se refletiu também em respostas de opinião e que evidenciavam preferências políticas.
Foram encontrados posicionamentos partidários que não condiziam com o conteúdo específico a ser avaliado.
Na questão 5, sobre o TNP, embora a própria questão revelasse a necessidade de se esclarecerem posicionamentos sobre o tema por países em posições diferentes no cenário internacional, houve posicionamentos parciais.
No entanto, foi a questão menos problemática na aplicação do padrão de resposta.
Por fim, ressaltam -se dois últimos comentários gerais sobre a prova.
Primeiro, foi revelada uma tendência de polarização e um gap, de um lado respostas muito fracas e, de outro, respostas muito boas.
Esse gap reflete uma situação cindida dos alunos que realizaram a prova de Relações Internacionais e foi identificado nas três questões, mesmo com variações entre elas.

% matched lemmas: cada, estar, perguntar
\end{textsample}
